\begin{abstract}
The GDPR's cooperation and consistency mechanism is a network of peer authorities collectively tasked with the uniform application of the law protecting the fundamental right to data protection. 
Since its introduction, the mechanism has been criticized, \emph{inter alia}, for affording the lead supervisory authority too large a role in setting the scope of investigations, leading to under-investigation and under-enforcement. 
The recently adopted Regulation 2025/2518 reforms this procedure by proceduralizing many aspects of the cooperation required between data protection authorities. This study examines the reformed flow in light of two of the sub-principles of good administration: the duty of care and the duty to give reasons. 
The reforms are assessed against network theories of peer pressure, asking whether the proceduralization offered by the Procedural Regulation will effectively curb the dominance of the lead supervisory authority and produce better enforcement outcomes, bolstering respect for the principle of good administration. 
While the reform brings improvements by providing formal steps for interaction, it does not address the structural issues of under-resourcing and diverging national priorities that have long characterized cooperation between data protection authorities. 
This may still leave concerned data protection authorities unable to meaningfully participate in the procedure and contribute to more careful and reasoned decisions. 
Under-enforcement of data protection law is especially troubling in connection with the ever-increasing scale of personal data processing and the speed of technological development, which require strong regulatory oversight to strike a balance between the protection of fundamental rights and the free movement of data. 
The success of this reform will depend on whether data protection authorities are willing and able to respect the strict procedural requirements it imposes and to fulfill their constitutional duty, under EU primary law, to protect the fundamental right to data protection.
\end{abstract}